\section{Plan de capacitación y gestión del cambio}

\subsection{Objetivo de la gestión del cambio}
\begin{indentedsubsection}
La estrategia de capacitación y gestión del cambio tiene como propósito fundamental facilitar y asegurar una transición tecnológica fluida para Perú Flores Travel Agency Tours Operador E.I.R.L. El objetivo es transitar de un modelo de gestión netamente operativo y manual ---altamente dependiente de libretas físicas, cálculos repetitivos en hojas de Excel y comunicaciones dispersas por WhatsApp--- hacia un entorno digital unificado, ágil y escalable en la nube.

Un pilar crítico de este plan es empoderar de manera directa a la Gerencia General (Silveria Flores) para que adquiera total autonomía tecnológica. Se busca que la gerencia no dependa exclusivamente de terceros o de su equipo de apoyo familiar para operar el sistema. El objetivo es que, al finalizar el proceso, la gerencia tenga el control absoluto de la herramienta, comprenda a la perfección la lógica del ecosistema Odoo y se sienta completamente segura al momento de gestionar su base de datos, elaborar propuestas comerciales y hacer seguimiento a sus operaciones diarias.
\end{indentedsubsection}

\subsection{Metodología de transferencia de conocimientos}
\begin{indentedsubsection}
Para garantizar una adopción exitosa y mitigar cualquier resistencia natural al cambio tecnológico, el plan descarta por completo el uso de manuales teóricos genéricos o sesiones magistrales. En su lugar, se implementa un enfoque pedagógico 100\% práctico y aplicado, denominado ``Aprender Haciendo''.

La capacitación está dirigida a la gerencia y a su equipo de apoyo familiar (su hija, quien colabora activamente en la agencia). Este entrenamiento se realizará utilizando los datos reales y los paquetes turísticos que la empresa ya comercializa (como los programas a Cusco, Arequipa y Puno), estructurándose en las siguientes áreas clave:

\begin{itemize}
    \item \textbf{A. Capacitación operativa en CRM y Base de Datos:} entrenamiento inmersivo para el ingreso y categorización de nuevos prospectos (leads) y la gestión del directorio de proveedores. Se enseñará a organizar a los clientes extranjeros dentro del embudo de ventas, eliminando definitivamente la necesidad de transcribir información desde cuadernos físicos a la computadora.
    \item \textbf{B. Dominio del Cotizador Automatizado:} instrucción práctica y detallada sobre cómo armar paquetes turísticos complejos utilizando el catálogo de servicios. La gerencia aprenderá a simular escenarios de rentabilidad, ajustando los porcentajes de ganancia (por ejemplo, 15\%, 20\% o 30\%) según el perfil del turista internacional. Asimismo, se practicará la generación de la propuesta comercial en formato PDF, dominando la configuración para ocultar el desglose del IGV por servicio, asegurando así una presentación limpia y profesional.
    \item \textbf{C. Control Financiero y Gestión de Cobros Fraccionados:} práctica exhaustiva en el registro de la trazabilidad financiera. Se capacitará al equipo en cómo asentar los depósitos y adelantos recibidos desde el extranjero (como los realizados vía Western Union) y cómo el sistema calcula automáticamente los saldos pendientes de pago, otorgando seguridad antes de la ejecución de los servicios turísticos.
    \item \textbf{D. Operación del Módulo de Encuestas (Feedback):} taller enfocado en el diseño y envío de los formularios de retroalimentación. La gerencia aprenderá a automatizar el envío de estas encuestas al finalizar el itinerario de los turistas, permitiendo consolidar testimonios reales que potenciarán el ``boca a boca'' digital y medirán la calidad de sus proveedores logísticos.
\end{itemize}
\end{indentedsubsection}

\subsection{Despliegue y retroalimentación}
\begin{indentedsubsection}
Para asegurar que el proceso de aprendizaje no interrumpa ni paralice las operaciones comerciales diarias de la agencia, el despliegue del plan de capacitación se ejecutará bajo un formato híbrido y sumamente flexible.

Las sesiones se programarán en estricta coordinación con la disponibilidad de la agenda de la gerencia, combinando jornadas de entrenamiento presencial en las instalaciones de la empresa con reuniones de revisión virtual. Las sesiones presenciales estarán enfocadas en la interacción directa con la plataforma, permitiendo a los usuarios resolver dudas en tiempo real mientras cotizan paquetes reales. Por su parte, las sesiones virtuales servirán para validar formatos, revisar avances estructurales (como el diseño del PDF) y afinar configuraciones previas al uso en vivo. Para mantener una comunicación fluida, soporte rápido y resolución de consultas operativas, se utilizará un canal oficial de coordinación mediante un grupo de WhatsApp dedicado exclusivamente al proyecto.
\end{indentedsubsection}

\subsection{Resultados esperados de la adopción tecnológica}
\begin{indentedsubsection}
De acuerdo con los objetivos estratégicos trazados en el Plan de Implementación, se espera que al culminar esta fase de capacitación, el 100\% del personal involucrado (tanto la gerencia como el equipo de apoyo) domine el ecosistema a nivel operativo.

El resultado cualitativo más trascendental será el drástico ahorro de tiempo y la completa eliminación de la dependencia de cálculos matemáticos manuales, sumas repetitivas o el uso del método de ``copiar y pegar'' en procesadores de texto. Al dotar a la Gerencia General de la confianza, velocidad y precisión necesarias para enviar cotizaciones profesionales directamente desde su equipo informático o dispositivo móvil, se garantiza la rentabilidad de la inversión tecnológica, asegurando que el sistema Odoo se convierta en el pilar central y definitivo para la escalabilidad a largo plazo de Perú Flores Travel Agency.
\end{indentedsubsection}
